\chapter{Carboxyhemoglobin}

\section*{What Is This Test?}
Carboxyhemoglobin measures the fraction of hemoglobin bound to carbon monoxide. It is the key laboratory measurement for suspected carbon-monoxide exposure and is usually obtained with a co-oximetry-capable blood-gas analyzer.

\section*{Alternative Names}
\begin{itemize}
  \item COHb
  \item Carbon Monoxide Level
  \item Carbon Monoxide Blood Test
  \item Co-Oximetry
\end{itemize}
\section*{Why This Test Is Ordered}
Testing is used after possible exposure to vehicle exhaust, faulty heating equipment, fires, industrial sources, or other carbon-monoxide environments, especially when headache, dizziness, confusion, chest symptoms, or unexplained neurologic illness occurs.

\section*{How This Test Is Performed}
Venous blood may be adequate for many diagnostic situations, while arterial blood may be collected when oxygenation and acid-base status also need evaluation. Co-oximetry directly measures oxyhemoglobin, deoxyhemoglobin, carboxyhemoglobin, and sometimes methemoglobin.

\section*{Preparation, Risks, and Limitations}
No fasting is required. Remove the patient from exposure and obtain the sample promptly because carboxyhemoglobin falls after breathing fresh air or receiving oxygen. A normal value after delay does not exclude a clinically important exposure. Smoking can produce a baseline elevation.

\section*{Understanding Results}
Higher values support carbon-monoxide exposure, but treatment decisions depend on symptoms, exposure history, pregnancy, neurologic findings, cardiac injury, and timing—not the number alone. Standard pulse oximetry cannot reliably distinguish oxyhemoglobin from carboxyhemoglobin.

\section*{References}
Centers for Disease Control and Prevention clinical guidance on carbon-monoxide poisoning.

\clearpage
\section*{Understanding Results and Follow-up}
The laboratory report should identify the specimen, method, units, reference interval or decision limit, and whether the result is positive, negative, abnormal, or inconclusive. Reference intervals vary with age, sex, pregnancy, specimen type, assay, and laboratory; a result outside a range is not automatically a diagnosis, and a result within a range does not always exclude disease. Discuss unexpected results with the ordering clinician, who can decide whether repeat or confirmatory testing, treatment, or referral is appropriate. Do not change medicines or delay urgent care based on an isolated result.


\section*{Regulatory Status and Authorization Date}
This chapter describes a test category, not a specific commercial product. FDA, CDSCO, EU IVDR, and MHRA approval or registration dates therefore cannot be assigned without the assay manufacturer, product name, instrument, and intended use. For laboratory-developed tests, an individual agency approval date may not exist. See the Preface for the interpretation of regulatory status.

\clearpage
